\documentclass[11pt]{article}

\usepackage[utf8]{inputenc}
\usepackage[T1]{fontenc}
\usepackage[margin=1in]{geometry}
\usepackage{amsmath, amssymb}
\usepackage{booktabs}
\usepackage{longtable}
\usepackage{array}
\usepackage{enumitem}
\usepackage{siunitx}
\usepackage{hyperref}
\usepackage{xcolor}
\usepackage{listings}

\hypersetup{
  colorlinks=true,
  linkcolor=blue!60!black,
  urlcolor=blue!60!black,
  citecolor=blue!60!black
}

\lstdefinestyle{repo}{
  basicstyle=\ttfamily\small,
  breaklines=true,
  frame=single,
  columns=fullflexible,
  keepspaces=true,
  showstringspaces=false
}

\title{Aircraft Trajectory Reconstruction: End-to-End Project Documentation}
\author{Repository Documentation Generated from the Current Codebase}
\date{\today}

\begin{document}

\maketitle

\begin{abstract}
This document provides a complete technical description of the aircraft trajectory reconstruction project implemented in this repository. It covers historical data ingestion from OpenSky, trajectory cleaning and deterministic train/test splitting, the great-circle and linear interpolation baselines, the constant-velocity Kalman filter and Rauch--Tung--Striebel smoother, the residual bidirectional LSTM gap-filling model, the training and evaluation pipelines, the emissions sensitivity analysis, and the FastAPI plus browser demo layer. The document is intentionally written as a single Overleaf-compatible \texttt{main.tex} file using standard packages only, so it can be uploaded and compiled directly on Overleaf with pdf\LaTeX{}.
\end{abstract}

\tableofcontents
\newpage

\section{Project Overview}
The repository implements an end-to-end workflow for reconstructing missing segments in aircraft trajectories. The main idea is to start from ADS-B/ADS-C flight tracks, remove noisy or irrelevant samples, train a model on complete trajectory windows, and then use the trained model to fill missing points in real or simulated gaps.

The codebase is organized around five major stages:
\begin{enumerate}[leftmargin=1.5em]
  \item Historical data ingestion from the OpenSky Trino database.
  \item Trajectory cleaning and deterministic train/test splitting into parquet files.
  \item Model development, including interpolation baselines, Kalman methods, and a residual bidirectional LSTM.
  \item Quantitative evaluation on held-out trajectories and a downstream emissions sensitivity study.
  \item A FastAPI service with a web demo for local and live OpenSky reconstructions.
\end{enumerate}

At the time of writing, the current workspace contains:
\begin{itemize}[leftmargin=1.5em]
  \item \textbf{7215} training-track parquet files under \texttt{data/clean/tracks}
  \item \textbf{1808} held-out test-track parquet files under \texttt{data/clean/test\_tracks}
  \item a checked-in held-out evaluation summary in \texttt{evaluation\_results.csv}
\end{itemize}

\section{Repository Layout}
Table~\ref{tab:repo-layout} summarizes the main project files.

\begin{longtable}{>{\raggedright\arraybackslash}p{0.28\textwidth} >{\raggedright\arraybackslash}p{0.64\textwidth}}
\caption{Core repository files and responsibilities}\label{tab:repo-layout}\\
\toprule
\textbf{File} & \textbf{Purpose} \\
\midrule
\endfirsthead
\toprule
\textbf{File} & \textbf{Purpose} \\
\midrule
\endhead
\texttt{data\_pipeline.py} & Queries OpenSky historical data via Trino, unnests track points, cleans trajectories, and writes deterministic train/test parquet splits. \\
\texttt{baseline.py} & Implements great-circle distance, spherical interpolation, altitude interpolation, and a generic gap-filling baseline. \\
\texttt{model.py} & Contains the constant-velocity Kalman filter, RTS smoother, residual BiLSTM model, and a wrapper fusion interface. \\
\texttt{train\_lstm.py} & Trains the residual BiLSTM, prepares windows, configures the TensorFlow runtime, and benchmarks validation/test gaps. \\
\texttt{evaluate.py} & Older qualitative evaluation script that hides one long gap and compares the great-circle baseline with Kalman methods. \\
\texttt{evaluate\_all.py} & Main held-out evaluation script used for report-style summary metrics over many random test gaps. \\
\texttt{phase6\_emissions.py} & Computes route-distance and CO\(_2\) proxy errors induced by gap reconstruction. \\
\texttt{api.py} & FastAPI service exposing health, local-flight listing, reconstruction, comparison demo, and OpenSky-backed live reconstruction. \\
\texttt{opensky\_client.py} & Small OpenSky REST client with OAuth client-credentials support and token caching. \\
\texttt{templates/index.html} & Browser UI for 2D/3D visualization, comparison, local reconstruction, and live OpenSky selection. \\
\texttt{weights/lstm\_residual\_trajectory.keras} & Main trained residual BiLSTM weights used by the service and evaluation scripts. \\
\texttt{evaluation\_results.csv} & Saved output from the official held-out evaluation script. \\
\bottomrule
\end{longtable}

\section{Environment and Dependencies}
The Python dependencies are intentionally lightweight except for the optional machine learning stack. The repository requires:
\begin{itemize}[leftmargin=1.5em]
  \item \texttt{requests}, \texttt{pandas}, \texttt{pyarrow}, \texttt{fastparquet}, and \texttt{python-dotenv} for data and I/O
  \item \texttt{numpy}, \texttt{scipy}, and \texttt{scikit-learn} for scientific computing
  \item \texttt{tensorflow} and \texttt{keras} for the LSTM model
  \item \texttt{fastapi} and \texttt{uvicorn} for the web service
\end{itemize}

The training script explicitly notes a platform caveat: native Windows TensorFlow versions newer than 2.11 do not use NVIDIA CUDA GPUs, so the code recommends WSL2 for GPU acceleration. When a GPU is available, the training script enables memory growth and can turn on mixed precision.

\section{Historical Data Pipeline}
\subsection{Source System}
The ingestion entrypoint is \texttt{data\_pipeline.py}. Historical trajectories are obtained from OpenSky's Trino-backed table \texttt{flights\_data4}. Each query is performed against a single OpenSky day partition, and the nested \texttt{track} array is unnested into point-level rows.

The default airport filter is: \texttt{EDDF}, \texttt{EGLL}, \texttt{LFPG}, \texttt{EHAM}, \texttt{LEMD}, \texttt{LIRF}, \texttt{LOWW}, \texttt{LSZH}, \texttt{EKCH}, and \texttt{EDDM}.

At a high level, the SQL performed by the pipeline is:
\begin{lstlisting}[style=repo]
SELECT f.icao24, f.firstseen, f.lastseen,
       f.estdepartureairport, f.estarrivalairport, f.callsign,
       tp_time, tp_lat, tp_lon, tp_alt, tp_heading, tp_onground
FROM flights_data4 f
CROSS JOIN UNNEST(f.track) AS t(
    tp_time, tp_lat, tp_lon, tp_alt, tp_heading, tp_onground
)
WHERE f.day = <partition>
  AND (
      f.estdepartureairport IN (<airport list>)
      OR f.estarrivalairport IN (<airport list>)
  );
\end{lstlisting}

\subsection{Daily Partitioning}
The helper \texttt{day\_partitions(days)} generates UTC-midnight timestamps for the previous \(N\) full UTC days, excluding the current day. This keeps data retrieval aligned with OpenSky's day partition structure and avoids partial-day inconsistencies.

\subsection{Track Cleaning}
Each grouped flight trajectory is processed by \texttt{clean\_track(df)}:
\begin{enumerate}[leftmargin=1.5em]
  \item sort by timestamp
  \item drop duplicate timestamps
  \item remove points with missing latitude, longitude, or barometric altitude
  \item interpret \texttt{on\_ground} as a nullable boolean and fill missing values with \texttt{False}
  \item remove points with altitude at or above \SI{15000}{\metre}
  \item remove points flagged as on-ground
\end{enumerate}

This cleaning stage is designed to keep airborne trajectory segments only and to remove sparse or corrupted records before model training.

\subsection{Deterministic Train/Test Split}
The pipeline does not randomly reshuffle flights every run. Instead, it hashes the string key \texttt{icao24|firstseen} using CRC32 and routes each flight to train or test deterministically. If \(h_f\) denotes that hash for flight \(f\), then:
\[
\text{test}(f)=
\begin{cases}
1, & \text{if } h_f \bmod 100 < 100r\\
0, & \text{otherwise}
\end{cases}
\]
where \(r\) is the configured test ratio.

This is an important design choice: the same flight always lands in the same split across repeated runs, so evaluation is reproducible at the flight level.

\subsection{Stored Output Format}
Each cleaned trajectory is saved to a parquet file named \texttt{<icao24>\_<firstseen>.parquet}
and contains at least the following columns:
\begin{itemize}[leftmargin=1.5em]
  \item \texttt{time}
  \item \texttt{latitude}
  \item \texttt{longitude}
  \item \texttt{baro\_altitude}
  \item \texttt{true\_track}
  \item \texttt{on\_ground}
  \item \texttt{icao24}
  \item \texttt{callsign}
  \item \texttt{flight\_start}
  \item \texttt{flight\_end}
\end{itemize}

\section{Geometric and Physical Foundations}
\subsection{Great-Circle Distance}
The repository consistently measures geographic point-to-point error with the haversine great-circle distance. For latitudes and longitudes in radians:
\begin{align}
\Delta \phi &= \phi_2 - \phi_1, \\
\Delta \lambda &= \lambda_2 - \lambda_1, \\
a &= \sin^2\!\left(\frac{\Delta \phi}{2}\right)
  + \cos(\phi_1)\cos(\phi_2)\sin^2\!\left(\frac{\Delta \lambda}{2}\right), \\
d_{\text{gc}} &= 2R \arctan2(\sqrt{a}, \sqrt{1-a}),
\end{align}
with Earth radius \(R = \SI{6371}{\kilo\metre}\).

This formula appears in both \texttt{baseline.py} and \texttt{model.py}, and it is the core error measure used throughout evaluation.

\subsection{Great-Circle Interpolation}
Given two spherical endpoints and an interpolation fraction \(f \in [0,1]\), the baseline computes an intermediate point using spherical linear interpolation:
\begin{align}
\delta &= d_{\text{gc}} / R, \\
A &= \frac{\sin((1-f)\delta)}{\sin(\delta)}, \\
B &= \frac{\sin(f\delta)}{\sin(\delta)}.
\end{align}
The corresponding Cartesian coordinates are:
\begin{align}
x &= A\cos\phi_1\cos\lambda_1 + B\cos\phi_2\cos\lambda_2, \\
y &= A\cos\phi_1\sin\lambda_1 + B\cos\phi_2\sin\lambda_2, \\
z &= A\sin\phi_1 + B\sin\phi_2.
\end{align}
The interpolated latitude and longitude are then:
\begin{align}
\phi &= \arctan2\!\left(z, \sqrt{x^2+y^2}\right), \\
\lambda &= \arctan2(y,x).
\end{align}

Altitude is interpolated linearly:
\[
\hat{h}(f) = h_1 + f(h_2-h_1).
\]

\subsection{Constant-Velocity Kalman Filter}
The filter state is:
\[
\mathbf{x}_t =
\begin{bmatrix}
\text{lat}_t & \text{lon}_t & \text{alt}_t & v^{\text{lat}}_t & v^{\text{lon}}_t & v^{\text{alt}}_t
\end{bmatrix}^{\top}.
\]

For time step \(\Delta t\), the transition matrix is:
\[
\mathbf{F}(\Delta t)=
\begin{bmatrix}
1&0&0&\Delta t&0&0\\
0&1&0&0&\Delta t&0\\
0&0&1&0&0&\Delta t\\
0&0&0&1&0&0\\
0&0&0&0&1&0\\
0&0&0&0&0&1
\end{bmatrix}.
\]

The filter uses standard predict/update equations:
\begin{align}
\mathbf{x}_{t|t-1} &= \mathbf{F}\mathbf{x}_{t-1|t-1}, \\
\mathbf{P}_{t|t-1} &= \mathbf{F}\mathbf{P}_{t-1|t-1}\mathbf{F}^{\top} + \mathbf{Q}, \\
\mathbf{y}_t &= \mathbf{z}_t - \mathbf{H}\mathbf{x}_{t|t-1}, \\
\mathbf{S}_t &= \mathbf{H}\mathbf{P}_{t|t-1}\mathbf{H}^{\top} + \mathbf{R}, \\
\mathbf{K}_t &= \mathbf{P}_{t|t-1}\mathbf{H}^{\top}\mathbf{S}_t^{-1}, \\
\mathbf{x}_{t|t} &= \mathbf{x}_{t|t-1} + \mathbf{K}_t\mathbf{y}_t, \\
\mathbf{P}_{t|t} &= (\mathbf{I} - \mathbf{K}_t\mathbf{H})\mathbf{P}_{t|t-1}.
\end{align}

Only position is measured directly, so \(\mathbf{H}\) extracts the position part of the state. Gap points are handled by prediction-only steps: if a point is unobserved, no measurement update is applied.

\subsection{RTS Smoother}
The smoother in \texttt{model.py} performs a forward Kalman pass followed by a backward Rauch--Tung--Striebel pass. If \(\mathbf{x}_{t|T}\) is the smoothed estimate and \(\mathbf{x}_{t|t}\) is the filtered estimate, the smoothing gain is:
\[
\mathbf{D}_t = \mathbf{P}_{t|t}\mathbf{F}^{\top}
\left(\mathbf{F}\mathbf{P}_{t|t}\mathbf{F}^{\top} + \mathbf{Q}\right)^{-1},
\]
and the smoothed state update is:
\[
\mathbf{x}_{t|T}
=
\mathbf{x}_{t|t}
\mathbf{D}_t
\left(
\mathbf{x}_{t+1|T} - \mathbf{x}_{t+1|t}
\right).
\]

This bidirectional processing is useful for gap reconstruction because it uses context from both sides of the missing segment.

\section{Baseline and Classical Reconstruction Methods}
\subsection{Great-Circle Baseline}
The simplest baseline in this project connects the last known point before a gap to the first known point after a gap by a great-circle arc. It is intentionally strong: unlike naive Euclidean interpolation in latitude/longitude, it respects the geometry of the sphere.

\subsection{Linear Latitude/Longitude Baseline}
For comparison, the official evaluation also includes a purely linear baseline in latitude/longitude space:
\[
\hat{\mathbf{p}}_j = \mathbf{p}_{\text{before}} + \alpha_j \left(\mathbf{p}_{\text{after}} - \mathbf{p}_{\text{before}}\right),
\quad
\alpha_j = \frac{j}{G+1}, \quad j=1,\dots,G,
\]
where \(G\) is the number of hidden gap points.

\subsection{Kalman Filter and Smoother}
The classical alternatives are exposed through \texttt{FusionTrajectoryModel}. These methods are mainly used in the older visual comparison script \texttt{evaluate.py}, which hides a \SI{20}{\minute} middle segment from a held-out flight and compares:
\begin{itemize}[leftmargin=1.5em]
  \item the great-circle baseline,
  \item the forward constant-velocity Kalman filter,
  \item the RTS smoother.
\end{itemize}

This script is best viewed as a qualitative or legacy comparison stage. The repository's main quantitative LSTM evaluation is \texttt{evaluate\_all.py}.

\section{Residual Bidirectional LSTM Model}
\subsection{Why a Residual Formulation?}
The BiLSTM does not predict absolute coordinates from scratch. Instead, it learns a residual correction on top of a linear interpolation baseline. This has two advantages:
\begin{enumerate}[leftmargin=1.5em]
  \item the network only needs to learn where simple interpolation is wrong
  \item the zero-initialized output layer makes the initial prediction equal to the baseline
\end{enumerate}

\subsection{Context Window}
The model expects:
\begin{itemize}[leftmargin=1.5em]
  \item \(C = 20\) points before the gap
  \item \(G = 10\) hidden gap points to reconstruct
  \item \(C = 20\) points after the gap
\end{itemize}
so the total context sequence length is:
\[
L = 2C = 40.
\]

\subsection{Feature Construction}
Let the final point before the gap be the anchor \((\phi_a,\lambda_a,h_a,t_a)\). For every context point, the feature vector is:
\[
\mathbf{f}_i =
\begin{bmatrix}
10(\phi_i-\phi_a) \\
10(\lambda_i-\lambda_a) \\
(h_i-h_a)/10000 \\
\Delta t_i/3600
\end{bmatrix}.
\]

The time delta is computed between adjacent points within the before segment and within the after segment; for the first after-gap point, the previous time is the anchor time.

In code, the corresponding constants are \texttt{POSITION\_SCALE = 10}, \texttt{ALTITUDE\_SCALE\_METERS = 10000}, and \texttt{TIME\_SCALE\_SECONDS = 3600}.

\subsection{Residual Targets}
The baseline gap trajectory is built by linearly interpolating between the last pre-gap point and first post-gap point in both position and altitude. Denote baseline position and altitude at hidden index \(j\) by \((\mathbf{b}^{p}_j, b^{h}_j)\), and denote the ground truth by \((\mathbf{g}^{p}_j, g^{h}_j)\). The residual target is:
\[
\mathbf{r}_j =
\begin{bmatrix}
10(g^{\phi}_j - b^{\phi}_j) \\
10(g^{\lambda}_j - b^{\lambda}_j) \\
(g^{h}_j - b^{h}_j)/10000
\end{bmatrix}.
\]

At inference time, the model predicts \(\hat{\mathbf{r}}_j\), and the final reconstruction is:
\[
\hat{\mathbf{g}}_j =
\begin{bmatrix}
b^{\phi}_j + \hat{r}^{\phi}_j/10 \\
b^{\lambda}_j + \hat{r}^{\lambda}_j/10 \\
b^{h}_j + 10000\,\hat{r}^{h}_j
\end{bmatrix}.
\]

\subsection{Network Architecture}
The Keras model created by \texttt{create\_lstm\_model()} is:
\begin{center}
\begin{tabular}{ll}
\toprule
\textbf{Stage} & \textbf{Shape / Setting} \\
\midrule
Input & \((40, 4)\) \\
BiLSTM & 64 units, return sequences \\
Dropout & 0.2 \\
BiLSTM & 32 units, final state only \\
Dropout & 0.2 \\
Dense & 64 units, ReLU \\
Dense & \(10 \times 3 = 30\) outputs, zero kernel and bias init \\
Reshape & \((10, 3)\) residual prediction \\
\bottomrule
\end{tabular}
\end{center}

The optimizer is Adam and the loss is Huber with \(\delta=1.0\). Mean absolute error is tracked as an auxiliary metric.

\section{Training Procedure}
\subsection{Trajectory Loading}
\texttt{train\_lstm.py} loads cleaned parquet trajectories from:
\begin{itemize}[leftmargin=1.5em]
  \item \texttt{data/clean/tracks} for training and validation
  \item \texttt{data/clean/test\_tracks} for held-out benchmarking
\end{itemize}

Any trajectory shorter than \(2C+G = 50\) points is skipped.

\subsection{Train/Validation/Test Split}
The split happens in two stages:
\begin{enumerate}[leftmargin=1.5em]
  \item \textbf{Pipeline stage:} deterministic train/test split at the flight level using CRC32 hashing.
  \item \textbf{Training stage:} the training-flight set is randomly split 80/20 into train and validation trajectories using NumPy permutation with seed 42.
\end{enumerate}

\subsection{Window Sampling}
For each eligible flight, sliding windows of length \(2C+G\) are extracted. The training script uses a default \texttt{window\_stride} of 4, so only every fourth valid window is used. This reduces overlap and keeps the number of windows manageable.

\subsection{Runtime Configuration}
The default hyperparameters are:
\begin{itemize}[leftmargin=1.5em]
  \item epochs: 50
  \item batch size: 128
  \item context length: 20
  \item gap length: 10
  \item window stride: 4
  \item early stopping patience: 6
  \item seed: 42
\end{itemize}

On GPU systems, the code can enable mixed precision and TensorFlow memory growth. Best model weights are restored via early stopping and saved to \texttt{weights/lstm\_residual\_trajectory.keras}.

\section{Evaluation Methodology}
\subsection{Official Held-Out Evaluation}
The main report-quality evaluation is \texttt{evaluate\_all.py}. It:
\begin{enumerate}[leftmargin=1.5em]
  \item loads the held-out flight set from \texttt{data/clean/test\_tracks}
  \item samples 3 random hidden windows per eligible flight with seed 42
  \item reconstructs each hidden gap with the trained LSTM
  \item compares against great-circle and linear baselines
  \item writes per-gap outputs to \texttt{evaluation\_results.csv}
\end{enumerate}

For each gap and each method, the script computes:
\begin{align}
\text{mean error} &= \frac{1}{G}\sum_{j=1}^{G} d_{\text{gc}}(\mathbf{g}_j, \hat{\mathbf{g}}_j), \\
\text{median error} &= \operatorname{median}\left(d_{\text{gc}}(\mathbf{g}_j, \hat{\mathbf{g}}_j)\right), \\
\text{p95 error} &= \operatorname{percentile}_{95}\left(d_{\text{gc}}(\mathbf{g}_j, \hat{\mathbf{g}}_j)\right), \\
\text{path error} &= \ell(\hat{\mathbf{g}}_{1:G}) - \ell(\mathbf{g}_{1:G}),
\end{align}
where \(\ell(\cdot)\) is the total great-circle path length along the hidden segment.

\subsection{Current Saved Held-Out Results}
The checked-in \texttt{evaluation\_results.csv} contains \textbf{5418} gap samples. Since the evaluation script uses 3 samples per eligible flight, this corresponds to \textbf{1806} eligible test flights and 2 skipped test flights in the current workspace.

Table~\ref{tab:heldout-results} summarizes the current saved results.

\begin{table}[h!]
\centering
\caption{Held-out reconstruction results from the checked-in \texttt{evaluation\_results.csv}}
\label{tab:heldout-results}
\begin{tabular}{lrrrrr}
\toprule
\textbf{Method} & \textbf{Mean} & \textbf{Median} & \textbf{Mean p95} & \textbf{Median $|$path err$|$} & \textbf{Wins} \\
& \textbf{(km)} & \textbf{(km)} & \textbf{(km)} & \textbf{(km)} & \textbf{(of 5418)} \\
\midrule
LSTM & 18.113 & 2.213 & 44.226 & 1.293 & 2725 \\
Great-circle & 22.485 & 2.939 & 40.702 & 1.330 & 1415 \\
Linear & 22.780 & 2.948 & 41.091 & 1.329 & 1278 \\
\bottomrule
\end{tabular}
\end{table}

Using the mean of per-gap mean error as the primary score, the current saved LSTM improves over the great-circle baseline by:
\[
\frac{22.485 - 18.113}{22.485}\times 100 \approx 19.44\%.
\]

An important nuance is visible in Table~\ref{tab:heldout-results}: the LSTM wins more often and lowers the average mean error substantially, but the average p95 error remains slightly worse than the great-circle baseline. This suggests the model improves typical cases while still producing a small number of hard outliers.

\subsection{Visual Comparison Script}
\texttt{evaluate.py} is a separate script that hides a \SI{20}{\minute} central segment from a held-out flight, reconstructs it using the great-circle baseline and Kalman methods, reports median/p95/max error, and saves a visualization as \texttt{phase5\_comparison.png}. This script is useful for figures and sanity checks, but it is not the primary aggregate benchmark for the LSTM.

\section{Phase 6: Distance and Emissions Sensitivity}
\texttt{phase6\_emissions.py} asks a downstream question: if the hidden trajectory is reconstructed incorrectly, how much route distance and emissions error does that induce?

The script uses a simple emissions proxy:
\[
\text{fuel rate} = 3.0 \ \text{kg fuel/km},
\qquad
\text{CO}_2 = 3.16 \ \text{kg CO}_2/\text{kg fuel}.
\]
Therefore:
\[
\text{CO}_2\text{ per km} = 3.0 \times 3.16 = 9.48 \ \text{kg CO}_2/\text{km}.
\]

For each flight, the script:
\begin{enumerate}[leftmargin=1.5em]
  \item selects the center window of length \(2C+G\)
  \item hides the middle \(G\) points
  \item reconstructs the gap with the LSTM, great-circle baseline, and linear baseline
  \item measures both the gap-only distance error and the full-flight distance error after replacing the hidden points
  \item converts distance error to CO\(_2\) proxy error
\end{enumerate}

For a distance error \(\Delta d\), the CO\(_2\) proxy error is:
\[
\Delta \text{CO}_2 = 9.48 \, \Delta d.
\]

This phase is useful because a reconstruction model can look good geometrically while still biasing downstream operational quantities such as route distance.

\section{Service Layer and Demo Application}
\subsection{FastAPI Service}
The application server is implemented in \texttt{api.py}. It loads the trained LSTM once at startup and falls back to the great-circle method when the model is unavailable or when a gap lacks enough context.

The core local reconstruction policy is:
\begin{enumerate}[leftmargin=1.5em]
  \item scan the time-sorted track
  \item treat any time jump larger than \texttt{GAP\_THRESHOLD\_SECONDS = 120} as a gap
  \item try the LSTM if at least 20 points exist on both sides
  \item otherwise fill the missing interval with 10 great-circle points
\end{enumerate}

\subsection{API Endpoints}
Table~\ref{tab:api-endpoints} summarizes the main routes.

\begin{table}[h!]
\centering
\caption{FastAPI endpoints}
\label{tab:api-endpoints}
\begin{tabular}{>{\ttfamily}p{0.38\textwidth} p{0.5\textwidth}}
\toprule
\textnormal{\textbf{Route}} & \textbf{Purpose} \\
\midrule
\textnormal{GET /health} & Service status, model status, and OpenSky auth status \\
\textnormal{GET /flights} & Lists locally stored flight IDs \\
\textnormal{GET /reconstruct/\{flight\_id\}} & Reconstructs a stored local parquet trajectory \\
\textnormal{POST /reconstruct} & Reconstructs a user-supplied raw trajectory \\
\textnormal{GET /demo/\{flight\_id\}?gap=n} & Hides a dense local segment and compares LSTM vs great-circle \\
\textnormal{GET /opensky/flights} & Lists live OpenSky flights in the current map bounds \\
\textnormal{GET /opensky/reconstruct/\{icao24\}?time=t} & Fetches a live OpenSky track and reconstructs it \\
\bottomrule
\end{tabular}
\end{table}

\subsection{OpenSky REST Integration}
\texttt{opensky\_client.py} provides a small REST client with OAuth client-credentials support. It reads:
\begin{itemize}[leftmargin=1.5em]
  \item \texttt{OPENSKY\_CLIENT\_ID}
  \item \texttt{OPENSKY\_CLIENT\_SECRET}
\end{itemize}
from the environment, requests access tokens, caches them with an expiry margin, and attaches them to REST calls.

The important design decision is architectural: the browser never talks directly to the OpenSky API. Instead, the browser talks to FastAPI, and FastAPI proxies the OpenSky requests. This avoids exposing API credentials in client-side JavaScript.

\subsection{Browser Demo}
The UI in \texttt{templates/index.html} is a single-page application built with:
\begin{itemize}[leftmargin=1.5em]
  \item Leaflet for the 2D map
  \item \texttt{globe.gl} / three-globe for 3D visualization
\end{itemize}

It provides two main tabs:
\begin{enumerate}[leftmargin=1.5em]
  \item \textbf{Reconstruct}: visualize actual gap filling on a local or live OpenSky track
  \item \textbf{Compare Methods}: hide a dense segment and compare LSTM vs great-circle against the truth
\end{enumerate}

Additional demo features include:
\begin{itemize}[leftmargin=1.5em]
  \item 2D map and 3D globe toggle
  \item altitude exaggeration in the globe view
  \item downloadable reconstructed JSON
  \item OpenSky live-flight loading based on current map bounds
  \item per-gap comparison browsing over up to 10 evenly spaced hidden windows
\end{itemize}

The 3D globe uses an altitude exaggeration formula:
\[
\text{globe altitude} =
\frac{\text{altitude\_m}}{1000}
\cdot
\frac{a_{\text{exag}}}{R_{\text{Earth}}},
\]
with \(R_{\text{Earth}} = 6371\) km and default \(a_{\text{exag}} = 40\).

\section{How to Run the Project}
\subsection{Data Ingestion}
\begin{lstlisting}[style=repo]
python data_pipeline.py --days 2 --clean --test-ratio 0.2
\end{lstlisting}

\subsection{Model Training}
\begin{lstlisting}[style=repo]
python train_lstm.py --epochs 50 --batch-size 128 --window-stride 4
\end{lstlisting}

\subsection{Official Held-Out Evaluation}
\begin{lstlisting}[style=repo]
python evaluate_all.py
\end{lstlisting}

\subsection{Kalman Comparison Figure}
\begin{lstlisting}[style=repo]
python evaluate.py
\end{lstlisting}

\subsection{Emissions Sensitivity}
\begin{lstlisting}[style=repo]
python phase6_emissions.py
\end{lstlisting}

\subsection{API and Demo}
\begin{lstlisting}[style=repo]
uvicorn api:app --reload
\end{lstlisting}

\section{Strengths, Caveats, and Extension Points}
\subsection{Strengths}
The repository has several strong engineering choices:
\begin{itemize}[leftmargin=1.5em]
  \item deterministic train/test splitting at the flight level
  \item a strong spherical baseline, not just naive straight-line interpolation
  \item a residual LSTM formulation that simplifies learning
  \item explicit held-out evaluation on flights the model never saw during training
  \item an integrated demo that exposes both local and live OpenSky trajectories
\end{itemize}

\subsection{Caveats}
Important caveats include:
\begin{itemize}[leftmargin=1.5em]
  \item the model assumes a fixed context length (20) and fixed gap length (10)
  \item the current LSTM improves the average mean error but still shows some heavy-tail outliers
  \item the altitude cap and on-ground filtering are hand-crafted cleaning rules
  \item native Windows TensorFlow may not use CUDA GPUs
  \item several UI strings in the HTML file still show encoding artifacts, although this does not affect the mathematical or service logic
\end{itemize}

\subsection{Natural Extensions}
Possible next extensions are:
\begin{itemize}[leftmargin=1.5em]
  \item variable-length gap prediction instead of fixed 10-point outputs
  \item richer features such as heading, speed, or airport priors
  \item probabilistic uncertainty estimates for reconstructed points
  \item explicit handling of long gaps with hierarchical or multi-scale models
  \item historical OpenSky flight search in the live demo, not just current state browsing
\end{itemize}

\section{Overleaf Usage}
This file is already structured as a single Overleaf-compatible entrypoint:
\begin{itemize}[leftmargin=1.5em]
  \item filename: \texttt{main.tex}
  \item no custom document class
  \item no shell-escape requirement
  \item no external bibliography file
  \item only standard packages commonly available in Overleaf
\end{itemize}

To use it on Overleaf:
\begin{enumerate}[leftmargin=1.5em]
  \item upload the repository (or at least \texttt{main.tex}) to Overleaf
  \item set \texttt{main.tex} as the main document if Overleaf does not choose it automatically
  \item compile with pdf\LaTeX{}
\end{enumerate}

\section{Conclusion}
This repository implements a complete aircraft trajectory reconstruction system, beginning with OpenSky historical ingestion and ending with an interactive browser demo. The core modeling idea is a residual bidirectional LSTM that learns corrections relative to a simple interpolation baseline, trained on fixed windows extracted from cleaned flight tracks. The project also keeps classical geometric and Kalman-based methods in the loop, which is valuable both for benchmarking and for understanding where the learned model helps most.

From a systems perspective, the project is notable because it does not stop at model training: it includes deterministic data management, reproducible evaluation, downstream emissions analysis, a service API, and a visualization interface that can operate on both local stored flights and live OpenSky data. That makes it not just a model prototype, but a full applied ML and geospatial reconstruction workflow.

\end{document}
